\documentclass[10pt,a4paper]{article}

\usepackage[T1]{fontenc}
\usepackage[utf8]{inputenc}
\usepackage{lmodern}
\usepackage{amsmath,amssymb,amsthm}
\usepackage{booktabs,array}
\usepackage{geometry}
\usepackage{xcolor}
\usepackage{enumitem}
\usepackage{microtype}
\usepackage{hyperref}
\usepackage{fancyhdr}

\geometry{margin=1in}
\hypersetup{
  colorlinks=true,
  linkcolor=blue!55!black,
  citecolor=blue!55!black,
  urlcolor=blue!55!black,
  pdftitle={Equivalent Signing-Key Recovery against D-James q5/128},
  pdfauthor={ePrint Signature Audit}
}
\setlist{nosep,leftmargin=*}
\setlength{\parindent}{1em}
\setlength{\parskip}{0.25em}
\pagestyle{fancy}
\fancyhf{}
\fancyhead[L]{Equivalent signer against D-James q5/128}
\fancyhead[R]{7 October 2026}
\fancyfoot[C]{\thepage}
\renewcommand{\headrulewidth}{0.3pt}

\newtheorem{theorem}{Theorem}
\newtheorem{lemma}{Lemma}
\newtheorem{observation}{Observation}
\newcommand{\F}{\mathbb{F}}
\newcommand{\K}{\mathbb{K}}
\newcommand{\Tr}{\operatorname{Tr}}
\newcommand{\rank}{\operatorname{rank}}
\newcommand{\Span}{\operatorname{span}}
\newcommand{\Gal}{\operatorname{Gal}}

\title{\bfseries Practical Equivalent Signing-Key Recovery against D-James q5/128\\
\large Hidden Matrix-Gabidulin Recovery and HFE-IP Completion}
\author{ePrint Signature Audit\\
\small AI-assisted cryptanalysis report}
\date{7 October 2026}

\begin{document}
\maketitle

\begin{abstract}
We give a practical public-key-only equivalent signing-key recovery against
the advertised q5/128 parameter row of D-James, IACR ePrint 2026/1650,
version~5.  The public Dragon signature--hash tensor is an unknown-basis
expansion of an $\F_5$-subcode of a two-dimensional generalized Gabidulin
code over $\F_{5^{94}}$.  We specialize Vin\c{c}otte's published structural
attack against Miranda to recover the hidden field expansion and projected
output trace support.  We then use the public signature--signature block to
decode the HFE coefficients as a dimension-three Gabidulin word with a
rank-four error, recover the two-dimensional internal-perturbation plane,
complete an equivalent projected HFE-IP map, and recover all Dragon linear
maps.

The attack uses only public coefficients.  On a sanitized full-parameter
public key, independent recovery took 97.62 seconds and 192.25 MB.  Completing
the equivalent key and signing took 6 minutes 27.45 seconds and 104,960 KiB
peak RSS.  The recovered signer produced a nonzero 94-symbol signature that
satisfies all 73 original public equations for a fresh message and the
paper's all-ones public target.  Direct public evaluation rejected the same
signature for every one of 256 salts of a changed message.  The attack makes
no signing queries and does not recover or use the nominal secret key.

No current version-5 implementation is available, and the paper leaves its
hash, salt, field, and wire formats unspecified.  Our full coefficient
instance uses the exact q5 algebra and distribution stated in version~5,
generated by a pinned detailed implementation that predates the nonzero
right-hand side.  We instantiate the new right-hand side with the paper's
all-ones example and use the pinned implementation's SHAKE256 transcript for
the fresh-message demonstration.  Equivalent-key recovery signs arbitrary
public hash vectors and is independent of those byte-level choices.  The
result was AI-discovered and independently AI-reproduced, including two
additional freshly generated public keys; independent human verification is
pending.
\end{abstract}

\noindent\fcolorbox{black!45}{black!4}{\parbox{0.94\textwidth}{
\textbf{Result status.} Practical full-parameter design attack: public-key-only
equivalent signing-key recovery and fresh signing for D-James q5/128.  The
component Gabidulin-recovery techniques are published work and are credited
below.  The exact-target identification, completion, and application were not
located previously.  Independent human verification is pending.}}

\section{Introduction}

D-James~\cite{djames} combines HFE with minus and internal-perturbation
modifiers and adds Dragon bilinear terms between the signature variables and
message-hash variables.  The Dragon construction permits a hash dimension
larger than the number of public equations while keeping the signature short.
Version~5 changes verification from a homogeneous equation to a fixed public
nonzero right-hand side, blocking the zero-signature forgery that affects the
earlier homogeneous formulation.

The same Dragon terms reveal a different structural weakness.  After the
hidden output map and minus projection, every public hash slice is a base-field
expansion of a word in a small generalized Gabidulin code.  D-James publishes
only a random $\F_q$-subcode, but recent work on matrix-Gabidulin masking shows
how to recover the hidden extension-field structure.  Le attacks Enhanced
Gabidulin Matrix Codes~\cite{le}; Vin\c{c}otte develops a particularly direct
recovery for the Miranda signature scheme~\cite{vincotte}.  Random
$\F_q$-subcodes and their residual invariants were also studied by Metouk\'e
et al.~\cite{metouke}.

Our contribution is to identify D-James's public Dragon tensor as an instance
of this hidden-code problem and to carry the recovery through every remaining
public block to a practical signer.  The cross tensor recovers the field basis
and output trace support but does not itself expose an HFE inverse.  The public
signature--signature block supplies the missing information.  After
polarization, its HFE component is a sparse linearized polynomial and its
internal-perturbation component is a low-rank error.  Rank-metric decoding and
base-field linear algebra recover an equivalent central map, including the IP
plane and all projected coefficients.  The resulting object signs fresh
digests against the fixed nonzero target.

Our concrete contributions are:
\begin{itemize}
  \item an exact trace-adjoint derivation showing that the q5 Dragon tensor is
  an $\F_5$-subcode of an expanded two-dimensional generalized Gabidulin code;
  \item a normalized bivariate determinantal recovery that avoids the large
  generic Macaulay matrix and recovers the full field basis at $n=94$;
  \item a public rank-metric decoder for the q5 HFE coefficients and
  rank-four IP residual, followed by equivalent IP and Dragon completion;
  \item a full q5/128 execution producing a fresh accepted signature with zero
  signing queries, plus public-only and changed-message controls; and
  \item a sanitized reproducer whose attacker inputs contain no deterministic
  key-generation seed or secret-derived validation data.
\end{itemize}

\paragraph{Attack-history statement.}
We located no previous public attack against exact D-James version~5 in
searches as of 7 October 2026.  A final search covered the exact scheme name,
title, ePrint identifier, authors, source repositories, issue and commit
records, and combinations with attack, cryptanalysis, forgery, Gabidulin,
matrix code, Miranda, equivalent key, and related terms.  We separately
searched the three structural papers used here.  None names D-James or
completes its Dragon, HFE-IP, minus, and fixed-target layers.  The negative
search has moderate confidence because the target and decisive structural
paper are recent.  We make no absolute priority claim.

\paragraph{Technique attribution.}
The hidden matrix-Gabidulin recovery is a specialization of Vin\c{c}otte's
attack on Miranda~\cite{vincotte}, which builds on Le's EGMC
cryptanalysis~\cite{le}; we claim no novelty for those primitives.  We also
claim no novelty for Gabidulin decoding, trace-adjoint representations, or
low-rank internal-perturbation analysis.  The result here is their
exact-target identification and composition into a public-only equivalent
signer for D-James.

\section{Exact target and scope}

\subsection{The q5/128 row}

The target is D-James in ePrint 2026/1650, version~5, dated 24 September
2026~\cite{djames}.  The audited PDF has SHA-256 digest
\begin{center}
\small\ttfamily
3cf3347c43d37a2b9189d3ab20c44ec12541bd945193c5604efc439224b46f58.
\end{center}
The affected advertised row is
\[
 (q,n,m,a,n_y,r,D)=(5,94,73,21,111,2,6),
\]
with HFE monomials $X^2$ and $X^6=X^{q+1}$.  Here $a=n-m$ is the number
of minus coordinates, $n_y$ is the number of public hash coordinates, $r$ is
the number of internal-perturbation forms, and $D$ is the central degree.
The paper assigns this row 128-bit classical security and a 219-bit fast
signature.

Let $P:\F_5^n\times\F_5^{n_y}\to\F_5^m$ be the public quadratic/bilinear
map.  Version~5 verification accepts a signature vector $a$ and public hash
vector $y$ exactly when
\[
                  P(a,y)=c,
\]
where $c$ is a fixed public nonzero parameter.  The paper gives the all-ones
vector as its concrete example.  Our execution uses $c=\boldsymbol{1}_m$.

\subsection{Available source boundary}

The paper-linked C/C++ repository was unavailable during the audit.  The
author's Sage notebook and the detailed independent Python/Rust implementation
both predate version~5 and use the earlier homogeneous relation.  We pin the
latter to commit
\begin{center}
\small\ttfamily 4328ab366c066554094033297eb1ef3ec8b21c6f.
\end{center}
It supplies a concrete q5 key generator whose coefficient distribution agrees
with the literal version-5 algebra: the two table-listed HFE monomials,
bilinear IP degree strictly below $q^d$ (hence $k<d$), two IP forms, Dragon
maps $L_0,L_1$, random secret input/output maps, and the 21-coordinate minus
projection.  We replace only the public right-hand side with the all-ones
example.

The pinned implementation's field polynomial, SHAKE256 hash-to-field map,
256-salt domain, root order, and encoding are implementation choices because
version~5 does not specify them.  Our algebraic recovery uses none of these
choices except the public coefficient representation.  The final transcript
uses the pinned SHAKE256 mapping to demonstrate a concrete fresh message.

\subsection{Adversarial input hygiene}

The initial deterministic experiment record included its generation seed as
metadata.  Although no attack script read that field, disclosure would permit
regenerating the nominal secret and therefore fails a strict public-only
review.  We removed all seed labels and seed digests, regenerated every
derived attacker input from a sanitized public-coefficient file, and reran
independent recovery from that file.  The sanitized public artifact has
SHA-256 digest
\begin{center}
\small\ttfamily
51c295dbf1ecd9ce77a890b5436391365f919505971f5a1b19e8f17341dc1e3a.
\end{center}
It contains only public coefficients, dimensions and orientation, the public
right-hand side, public messages and their public hash vectors, and source
provenance.  Secret-derived benchmark fixtures are not part of the attack or
public reproducer.

\section{The Dragon tensor is a hidden Gabidulin subcode}

Set $\K=\F_{q^n}$ and choose an unknown $\F_q$-basis
$\sigma=(\sigma_1,\ldots,\sigma_n)$ so that
\[
                  X=\sum_{i=1}^n a_i\sigma_i.
\]
Every retained output coordinate is an $\F_q$-linear functional on $\K$.
By nondegeneracy of the trace pairing, it has a unique representation
\[
                  \phi_j(Z)=\Tr_{\K/\F_q}(t_jZ).
\]
The minus map retains $m$ independent rows of an invertible output map, so
$t_1,\ldots,t_m$ are linearly independent over $\F_q$.

For q5, the Dragon contribution associated with hash coordinate $s$ is
\[
             \Lambda_{0,s}X+\Lambda_{1,s}X^q.
\]
Let $\gamma=(\gamma_1,\ldots,\gamma_n)$ be trace-dual to $\sigma$.  If
$c_{s,i,j}\in\F_q$ is the public coefficient of $a_i y_s$ in output
equation $j$, define the public slice after unexpansion by
\[
                  d_{s,j}=\sum_{i=1}^n\gamma_i c_{s,i,j}.
\]
Using $a_i=\Tr(\gamma_iX)$ and the trace-adjoint identity gives
\[
 d_{s,j}=t_j\Lambda_{0,s}+
          (t_j\Lambda_{1,s})^{q^{-1}}.
\]

\begin{lemma}
For a Dragon central map of q-degree $d$, all unexpanded public hash slices
belong to the generalized Gabidulin parent
\[
  \Span_{\K}\{t,t^{q^{-1}},\ldots,t^{q^{-d}}\}\subseteq\K^m.
\]
The public minus map shortens $t$ but does not destroy the parent as long as
the retained entries remain $\F_q$-independent.
\end{lemma}

For the target row, $d=1$.  The 111 public slices are therefore a random
$\F_5$-subcode of
\[
                  G_2(t)=\Span_{\K}\{t,t^{q^{-1}}\}.
\]
The parent has base-field dimension $2n=188$ and the public subcode has
codimension $188-111=77$.  Secret input and output maps change the expansion
basis and evaluation vector; they preserve this statement.

\section{Normalized public basis recovery}

Vin\c{c}otte's general recovery reduces the hidden-basis problem to a
rank-one extension-field system.  The special parent dimension two permits a
smaller normalized elimination.  Choose three public output columns.  From
the three entries of each compressed public matrix-code word, form public
matrices $A_0,A_1,A_2\in\F_5^{111\times94}$ such that
\[
                 A(u,v)=A_0+uA_1+vA_2.
\]
On the projective chart whose first compressed coordinate is nonzero, the
rank-one equations are equivalent to
\[
                 A(u,v)\gamma^T=0.
\]
Thus the hidden basis occurs on the determinantal locus
$\rank A(u,v)<94$.

We use three $94\times94$ maximal minors with row sets
\[
  1{:}94,\qquad 18{:}111,\qquad (1{:}93)\cup\{95\}.
\]
Computing these determinants symbolically was unnecessarily slow.  Each
minor has total degree at most 94, so we evaluate it on a $95\times95$ grid
over $\F_{5^3}$, interpolate in each variable, verify that the coefficients
descend to $\F_5$, and verify the degree bound.  The three minors contain
3,704, 3,635, and 3,625 nonzero terms.

Let $D_1,D_2,D_3$ denote the interpolated minors.  We compute
\[
 R_{12}=\operatorname{Res}_u(D_1,D_2),\qquad
 R_{13}=\operatorname{Res}_u(D_1,D_3).
\]
Both resultants have degree 8,836 in $v$.  Their gcd has degree 95 and
factorization degrees $1$ and $94$.  Adjoining a root of the unique degree-94
factor defines a recovered copy of $\F_{5^{94}}$; the gcd of all three minors
in $u$ is linear.  The resulting full pencil has rank 93 and a one-dimensional
right kernel.  Its kernel entries have base-field rank 94 and therefore form
a valid hidden expansion basis.

Unexpanding all 111 public slices gives extension-field rank two.  If $G$ is
their $\K$-row space, then
\[
                    \dim_{\K}(G\cap G^{[q]})=1.
\]
This intersection recovers $\Span_{\K}(t)$.  The recovered representative
$t$ has base-field rank 73, as required by the retained output map.

\section{Recovering the HFE-IP block}

The cross tensor supplies the hidden input basis and output trace support but
not yet a central inverse.  We recover the remaining structure from the
public $aa$ block.

For output equation $j$, let $Q_j$ be its public scalar quadratic form in the
signature variables.  Polarize it and define the trace-adjoint operator
$L_j:\K\to\K$ by
\[
 Q_j(X+Y)-Q_j(X)-Q_j(Y)=\Tr(YL_j(X)).
\]
All coefficients of $L_j$ are computable from the public form and the
recovered basis.

The q5 HFE core is
\[
                  \lambda_0X^2+\lambda_1X^{q+1}.
\]
Its contribution to the adjoint operator is
\[
 H_j(X)=2t_j\lambda_0X+t_j\lambda_1X^q+
        (t_j\lambda_1X)^{q^{-1}}.
\]
It therefore has linearized support $\{-1,0,1\}$, or consecutive support of
dimension three after a Frobenius shift.

Write the two IP forms as $z_\ell(X)=\Tr(\zeta_\ell X)$.  The remaining
operator $E_j=L_j-H_j$ has image contained in
\[
 \Span_{\F_q}\{\zeta_1,\zeta_2,t_j\mu_1,t_j\mu_2\},
\]
and hence rank at most four.  Each known $L_j$ is therefore a sparse
dimension-three Gabidulin word plus a rank-four error.

The implementation recovers an image-space annihilator
\[
 V(Y)=v_0Y+v_1Y^q+v_2Y^{q^2}+v_3Y^{q^3}+Y^{q^4}
\]
from the cyclic recurrence on the known linearized coefficients.  It then
solves a base-field linear system for the three erased error coefficients.
All 73 residual operators have rank exactly four.  Their common image-space
intersection is
\[
                 Z=\Span_{\F_q}\{\zeta_1,\zeta_2\},
                 \qquad \dim_{\F_q}Z=2.
\]

Choose $X_1,X_2$ trace-dual to any recovered basis of $Z$.  Modulo $Z$,
\[
                      E_j(X_\ell)=t_j\mu_\ell\pmod Z.
\]
Solving these membership constraints for all 73 outputs yields compatible
equivalent $\mu_1,\mu_2$.  The remaining image-in-$Z$ terms give the three
projected coefficients of the quadratic polynomial in $(z_1,z_2)$.  Each is
lifted through the surjective trace map
\[
 C\longmapsto (\Tr(t_1C),\ldots,\Tr(t_mC)).
\]
The 21 omitted output coordinates create equivalent choices, not an obstacle.
Eight independent random public evaluations confirm exact reconstruction of
the original projected $aa$ map.

\section{Dragon completion and fresh signing}

For each public hash coordinate $s$, the recovered cross slice obeys
\[
 d_{s,j}=t_j\Lambda_{0,s}+
         t_j^{q^{-1}}\Lambda_{1,s}^{q^{-1}}.
\]
Two output positions with an invertible $2\times2$ coefficient matrix recover
$\Lambda_{0,s},\Lambda_{1,s}$.  All 111 pairs fit all 73 public outputs.

For the fixed all-ones right-hand side, choose any lift $C\in\K$ satisfying
\[
                  \Tr(t_jC)=1\quad(1\le j\le73).
\]
The trace map has rank 73, so such a lift always exists.  More generally, the
same step lifts any public $m$-coordinate right-hand side.

Given a public hash vector $y$, the equivalent signer computes the recovered
Dragon coefficients, enumerates all 25 values of $(z_1,z_2)\in\F_5^2$, and
roots the resulting degree-six polynomial over $\K$.  It retains a root only
if the two recovered IP forms equal the guessed values, then returns
\[
                    a_i=\Tr(\gamma_iX).
\]

The full execution found a valid root after 19 degree-six attempts at salt 0
for the previously unsigned message
\begin{center}
\ttfamily D-James Miranda transfer fresh message.
\end{center}
The output is a nonzero vector of 94 elements of $\F_5$.  Direct evaluation
of the original 4,465 public $aa$ vectors and 111 public cross slices returns
$\boldsymbol{1}_{73}$.  Reusing the same vector for the changed message
\begin{center}
\ttfamily D-James Miranda transfer changed message
\end{center}
fails for every salt in the pinned 256-salt domain.

\begin{theorem}
For the executed D-James q5/128 public key, the recovered public data form an
equivalent signing witness: they reconstruct the entire projected public map
and produce a fresh nonzero signature accepted by the ordinary public
verification equations for the fixed all-ones target.
\end{theorem}

\section{Implementation and measurements}

All algebra was exact.  Public export and verification use Python; field and
polynomial recovery use Magma V2.29.  Every reported attack run used one
worker on an AMD EPYC Processor under KVM.  Table~\ref{tab:measurements}
records the primary measurements.

\begin{table}[ht]
\centering
\small
\begin{tabular}{@{}p{0.31\linewidth}p{0.34\linewidth}p{0.27\linewidth}@{}}
\toprule
Stage & Certificate & Measured cost\\
\midrule
Three determinant interpolations & 3,704 / 3,635 / 3,625 terms
  & 19.51 / 19.57 / 19.61 s\\
Two resultants & degree 8,836 each & 5.35 / 5.31 s\\
Independent basis recovery & rank $93$, $\gamma$ rank $94$, parent rank $2$,
intersection $1$, $t$ rank $73$ & 97.62 s, 192.25 MB\\
HFE rank decode & support $\{0,1,93\}$; 73 error ranks equal 4
  & 1.52 CPU s after conversion\\
Post-recovery completion/signing & all 111 Dragon pairs; full public-map
certificate; fresh signature & 386.92 CPU s; 6:27.45 wall\\
Post-recovery peak RSS & --- & 104,960 KiB\\
Dependency-free public verification & 73/73 fresh equations; 0/256 changed
salts & 0.52 s, 37,376 KiB\\
Fresh-key executions & 2/2 additional keys; 73/73 fresh equations and 0/256
changed salts in each trial & 98--99 s recovery plus 391--392 s signing\\
\bottomrule
\end{tabular}
\caption{Primary q5/128 measurements.}
\label{tab:measurements}
\end{table}

The generic transferred Miranda model has parent dimension $\kappa=2$,
subcode codimension $\ell_s=77$, and $\rho=n+\ell_s=171$.  Its degree-three
Macaulay column count is
\[
                 \binom{\rho+2}{3}=848{,}046.
\]
A dense base-field realization was estimated near 194 GiB.  The normalized
two-variable determinant route stays below 0.21 GiB and finishes in minutes.

The dependency-free verifier independently recomputes all 512 SHAKE256 hash
vectors rather than trusting those stored with the public key.  It then
evaluates the public equations directly from the sanitized coefficient JSON
and claimed signature.  A separate Magma recovery regenerated every
determinant from the sanitized JSON and loaded no prior determinant,
resultant, field modulus, recovered basis, output trace support, key seed, or
secret fixture.

\paragraph{Fresh-key reproduction.}
Two additional full executions started from independently generated,
sanitized public keys and imported no primary-key basis or recovery
checkpoint.  Trial~A used the public-input digest
\begin{center}
\small\ttfamily
68a358ce0c20fbe7788319d379a6038f79421875a6c690837ccb771aef168401.
\end{center}
Its clean recovery took 98.329 seconds and 224.28 MB.  Completion and signing
took 391.33 CPU seconds (391.871 seconds wall), found salt~0 after 14 root
attempts, and produced a nonzero signature accepted by all 73 equations.  An
independent direct verifier recomputed all 512 hash vectors and found no
accepting changed-message salt out of 256.

Trial~B used the public-input digest
\begin{center}
\small\ttfamily
96810a4a5548a6a04aa8e176df12b51e137a6588e475293c5a8d68080ea1da47.
\end{center}
Its three determinant interpolations took 19.55, 19.56, and 19.65 seconds;
the resultant stage took 39.37 seconds with 217,088 KiB peak RSS.  Completion
and signing took 390.80 CPU seconds (391.36 seconds wall), found salt~2 after
63 root attempts, and again gave 73/73 acceptance with zero accepting
changed-message salts out of 256.  The accompanying
\texttt{fresh-key-evidence.json} pins the full public-input, recovered-basis,
attack-transcript, forgery, and verification hashes for both trials.

Together with the primary execution, all three tested keys succeeded.  This
small sample demonstrates transfer across independently generated keys; it
does not support a precise statistical success-rate estimate.

\section{Limitations and claim boundary}

\paragraph{Version and wire format.}
No conforming version-5 implementation was available.  We claim the literal
paper algebra and its q5/128 coefficient distribution, not conformance to an
unpublished byte format.  The attack reconstructs an equivalent signer for
arbitrary public hash vectors; a future conforming implementation can supply
its own hash vector to the same recovered map.

\paragraph{IP range.}
Version~5 says that the bilinear IP polynomial has degree strictly less than
$q^d$, giving $k<d$ and matching the executed instance.  An older author
notebook uses $k\le d$.  We do not include that alternate branch in the
full-parameter claim.

\paragraph{Other parameter rows.}
The q5/256 row has the same two-variable parent structure at a larger
extension degree, but we did not need to execute it to break an advertised
security row.  The q2 and q4 rows have parent dimension three and require a
multivariate recovery beyond this bivariate implementation.  The q13 and q23
rows were not measured.  We claim only q5/128.

\paragraph{Equivalent rather than nominal key.}
Frobenius and scalar equivalences and the output kernel mean that the recovered
coefficients need not equal the nominal secret values.  This distinction does
not reduce impact: the recovered object reconstructs the public map and signs
new digests.

\paragraph{Multi-key evidence.}
The primary deterministic execution and two independently generated fresh-key
executions all succeeded.  Generality follows from the algebraic derivation;
the observed empirical count is only three successes in three executions, and
we do not infer a larger statistical success rate from that sample.

\section{Conclusion}

D-James's Dragon terms expose a small hidden generalized Gabidulin parent.
Published matrix-code recovery techniques reveal its field expansion and
output trace support; the public HFE-IP block can then be decoded and completed
to an equivalent central map.  The complete public-only attack breaks the
advertised q5/128 row in minutes and produces a fresh signature against the
fixed nonzero target.  The nonzero right-hand side prevents the earlier zero
forgery but does not protect the structural signing trapdoor.

\paragraph{AI disclosure.}
Target analysis, history search, derivation, implementation, adversarial
review, documentation, and independent reproduction were AI-assisted.
Independent AI reproduction is recorded as such.  Independent human
verification is pending.

\begin{thebibliography}{9}

\bibitem{djames}
J.~Patarin and A.~Roullet.
\newblock D-James: Ultra Short Multivariate Signatures.
\newblock IACR Cryptology ePrint Archive, Report 2026/1650, version 5, 2026.
\newblock \url{https://eprint.iacr.org/2026/1650}.

\bibitem{vincotte}
A.~Vin\c{c}otte.
\newblock How to break the Miranda signature scheme over matrix Gabidulin
codes.
\newblock arXiv:2609.30925, version 2, 2026.
\newblock \url{https://arxiv.org/abs/2609.30925}.

\bibitem{le}
T.~H.~Le.
\newblock Breaking ACDGV MinRank Gabidulin encryption schemes over matrix
codes.
\newblock arXiv:2608.03328, version 2, 2026.
\newblock \url{https://arxiv.org/abs/2608.03328}.

\bibitem{metouke}
F.~L.~Metouk\'e, H.~T.~Kalachi, H.~T.~Kamche, O.~Ndiaye, and S.~Ndjeya.
\newblock Subcodes of Lambda-Gabidulin Codes for Compact-Ciphertext
Cryptography.
\newblock arXiv:2604.18282, version 1, 2026.
\newblock \url{https://arxiv.org/abs/2604.18282}.

\end{thebibliography}

\end{document}
